\chapter{Terrain: Themes, Materials and the Slope Axis}
\label{ch:terrain}

\section{A Second Pass over a Finished World}

Painting terrain is not part of building it. It is a second pass over the realized
world, run after every structure has been stamped, and it adds no geometry at all.
It reads the terrain the world builder placed and rewrites its surface.

\begin{quote}
Where the stampers write rooms, cubes and objectives onto the terrain, this pass
dresses the terrain \emph{itself}, bucket by bucket, through whichever theme
applies; unthemed, every bucket stays the raw stone it already is.
\end{quote}

One invariant makes the whole model safe to state simply: \textbf{the painter
touches only stone}. Terrain arrives as a bedrock floor at $y=0$ with a stone
column above it, and every structure the stampers add is bedrock, wool, obsidian,
anything but stone. So the pass runs last and rewrites stone blocks only. Bedrock
is excluded by construction rather than by a special case.

Two refinements of that invariant were each learned as a defect. The test is over
the whole block and not over its id, because stone's id is shared by granite,
diorite, andesite and their polished forms; a plinth in polished diorite under a
red prop came back red until the check compared the pair rather than the id. And
the ground a room stands on is not a stamp, so it is painted like any other ground:
a room levels its own footprint by filling upward in stone, and the result comes
out in whatever the board's theme paints there.

\section{Four Buckets}

Every column splits into at most three painted parts and one kept part. The four
names are the shared vocabulary that every later stage speaks.

\begin{description}[leftmargin=1.6em,style=nextline,itemsep=4pt]
  \item[Rim] The edge lip: a single block, the top of an edge column. Not a volume.
    A map's rim is therefore a one-block-wide line tracing every lip, quartz
    against the grass.
  \item[Wall] The exposed riser below the rim, painted only on the face that is
    actually exposed. Stone below the shallowest drop stays stone, because it is
    buried behind the neighbour that made the drop and painting it would tint
    blocks nobody can see.
  \item[Surface] The top of the interior, as a layered stack to a configured depth
    rather than a fixed single block. One block of stone by default; grass over two
    dirt, three deep, in a meadow finish.
  \item[Fill] The interior body below the surface. This is the \emph{required}
    bucket: it claims whatever no other enabled bucket took, so a theme is never
    partial.
\end{description}

A freestanding nine-tall column against the void therefore reads, bottom to top,
one bedrock, seven wall, one rim.

\subsection{Edges over Eight Neighbours}

A column is an edge when any of its \emph{eight} neighbours is void or carries a
strictly lower surface. The four-neighbour test is wrong here, and the reason is
geometric: at a reentrant corner of an L or a comb, the corner block's only lower
neighbour is diagonal, so a four-neighbour test leaves a one-block hole and the rim
reads as broken. The extra four neighbours are exactly the reentrant corner cells,
so the rule thickens nothing and only completes the turn.

How wide a net the lip casts is itself authorable, through three nested tests.
\id{drop}, the default, lips only drops. \id{boundary} traces the whole plateau
outline, adding the edges that face something \emph{taller}, such as a room or an
approach wall. \id{void} narrows instead of widening: a column is rim only where a
neighbour is off the footprint entirely.

The difference matters on stepped ground. A staircase of five shapes each a block
above the last is a drop at every tread, so \id{drop} draws a lip on each one and
the body reads as five plateaus that happen to touch. \id{void} caps the outside
and leaves every tread bare, so it reads as one body with a rim around it.

\section{Materials and the Band Stack}

A material is either simple or composed. The composed one is the interesting case:

\begin{quote}
\id{LayeredMaterial} is a \id{BandStack}, an ordered run of \id{(material,
thickness)}, plus the \textbf{axis} it is read along. The stack states its bands
and what happens where they run out, and deliberately not the distance, because
the distance belongs to whoever is measuring it.
\end{quote}

That last clause is the design decision the whole section rests on. There is one
type with the axis named, not four types, because the bands, the thicknesses and
the run-out rule are identical whichever way the stack is read and only the
distance differs.

\begin{table}[htbp]
\caption{The four band axes. The stack is the same object in each row; only what
is measured changes.}
\label{tab:axes}
\small
\begin{tabularx}{\textwidth}{L{2.0cm} L{4.2cm} Y}
\toprule
\textbf{Axis} & \textbf{The distance} & \textbf{What it draws} \\
\midrule
\id{depth} & courses down from the top of the bucket (the default) &
  grass over two dirt; a wall's banded riser \\
\id{inward} & steps in from the landmass's void-facing edge &
  a cobble rim, two rings of stone brick, then a field \\
\id{height} & courses up from a stated world $Y$ &
  a hull's waterline; a tower banded by course, whatever the ground beneath it
  does \\
\id{slope} & how steeply the ground is inclined here, 0 to 89 degrees from level &
  meadow on the flat, coarse dirt on the shoulder, bare rock on the face of the
  same hill \\
\bottomrule
\end{tabularx}
\end{table}

\subsection{Running Out}

\id{ending} decides what happens past the last band, and \id{beyond} says what
shows there. Under \id{repeat} the last band claims everything further along, which
is right where the stack owns its whole space. A ring stack usually does not own
its whole space: it is meant to run a few rings in and then stop, so it ends
\id{handOver}, and \id{beyond} answers for everything past it.

This is the mechanism behind a trap worth stating in its own right, because it
produces documents that look wrong and are not. A stack does not cycle. Striping a
face sixteen courses tall with a four-course pattern under \id{repeat} means writing
the four courses out four times, because past the last band the last band claims
everything.

\subsection{Why the Height Axis Exists}

\id{depth} and \id{inward} both measure from the shape the block is in, so a colour
change along either splits nothing. A colour change stated as \emph{geometry} is a
different matter: a layer holds one theme, so banding a form by splitting it into a
layer per colour costs a layer per band, and that, rather than the shape, is what a
sculpture costs.

Measured over nine models, the layer count a shape alone needs against the count it
actually took: a starship 2 against 4, a robot 5 against 16, and a Rubik's cube,
which is a solid box with one run per column, 1 against 7. Read along \id{height}
the same banding is one material on one layer.

\section{The Slope Axis}

This is the load-bearing one, and it is the axis a board should be finished along.

\begin{quote}
The slope axis is an angle mask, and a thickness on it is a span of degrees.
\end{quote}

A stack of grass at 30, coarse dirt at 15 and cobblestone at 45 is therefore meadow
up to 29 degrees, coarse dirt from 30 to 44, and bare rock above.

\begin{lstlisting}[language=json]
{"kind":"layered","axis":"slope","stack":{"ending":"repeat","bands":[
   {"material":{"kind":"solid","id":2,"data":0},"thickness":30},
   {"material":{"kind":"solid","id":3,"data":1},"thickness":15},
   {"material":{"kind":"solid","id":4,"data":0},"thickness":45}]}}
\end{lstlisting}

The angle comes from Horn's 3$\times$3 gradient over the neighbouring surface tops,
so a ramp climbing one block a cell answers 45 exactly and level ground answers 0.
It is a fact about the \emph{surface}, carried by the whole column, which is why a
face's fill and the block on top of it resolve to the same band and a hillside comes
out finished top to bottom.

\begin{ruling}{Nothing Else Distinguishes a Hillside from a Field}
The wall bucket is the exposed riser, and a surface climbing one block a cell has
no riser at all, because its drop is one course and the top course already claims
it. Every other axis therefore paints a 45-degree hillside and a flat meadow alike.
A board finished on plan pieces or on height bands is painted flat from above
however much relief is under it.
\end{ruling}

The void is deliberately not a slope. A neighbour off the footprint, or one
carrying a structure, reads as level with the cell itself, so a coastline is flat
ground that stops, and the ground beside a building does not take the building's
roof for a gradient.

\subsection{The Window Is Two Cells}

The gradient is read two cells either side rather than one, and the reason is that
a gentle slope in a voxel world is a staircase.

\begin{quote}
Ground falling one block every four cells is four flat treads and a one-block
riser, and a window one cell wide reads each riser on its own, 27 degrees on the
riser and nothing on the treads, so a uniform gentle grade comes out speckled with
whatever the theme paints at 27.
\end{quote}

The tell is in the distribution rather than in the picture: read one cell either
side, a board authored at a one-block step puts a quarter of its ground in a spike
at 20 to 29 degrees with almost nothing at 30 to 39, which is a board reporting its
own quantum rather than its shape.

The arithmetic for why two is the right number is worth recording. A
\emph{sustained} slope reads its true angle at every window, so one block per four
cells is 14 degrees at windows of 1, 2, 3 and 4, and widening costs a real slope
nothing. What widening costs is a face, whose drop is spread over more cells: the
lip of a six-block drop reads 72 degrees at a window of one, 56 at two and 45 at
three, where it can no longer be told from a walkable ramp. The isolated one-block
contour that should be suppressed falls from 27 degrees to 14 to 9.5 over the same
range. Two is where the contour is already under any threshold an author would set
and the face is still unmistakably a face.

\subsection{Where the Cuts Go Is Read, Not Guessed}

\route{GET /map/\{slug\}/incline} answers how much ground stands in each ten
degrees. It is the only thing that says whether a band cut lands where the author
thinks it does, and \id{?window=} answers the same board read wider, so both
decisions a mask needs are made without an export.

Two corpus episodes show what happens without it. A first cut at 24 degrees put an
entire hill onto one band, and the board came out grey with green only at its
edges; recut at 32 after reading the histogram, it was right. On another board the
reading answered 77\% of cells under ten degrees, 22.6\% in the teens, 0.5\% above
twenty and nothing past forty, which is to say that a board about to be finished by
angle had almost no angle in it.

\section{Composition, and the Commonest Mistake}

The axes compose, and this removes the need for several knobs that would otherwise
be necessary. A surface material that is a \id{depth} stack whose first band is one
course thick, and is \emph{itself} an \id{inward} stack, rings the top course and
leaves the courses beneath it alone.

That same composition is what keeps a grass ring one block thick rather than three.
A palette that repeats grass down every course is, by the repository's own account,
the most-repeated authoring mistake in it.
